\documentclass{tufte-handout}
\usepackage[utf8]{inputenc}
\usepackage[T1]{fontenc}

\title{Propuesta de actividades}

\author{Angélica Pena,
        Raymundo Meza,
        Abraham Olivetti
       }
\date{}

\begin{document}

\maketitle


\textbf{Actividades sugeridas}

Me parece que para la primera parte del curso es buena idea que tengan claros algunos de los conceptos que marcan los debates clásicos de la filosofía del lenguaje, que además son coneptos que se recuperan en la literatura contemporánea. Es por estas razones que sugiero que una de las actividades continúas sea crear un glosario con algunas de las palabras clave de los artículos de Frege, Russell y Strawson. Otra sugerencia es reconstruir el argumento de Russel l que está en la página 51 del manual elegido para el curso y leer y discutir en clase las pp. 57 a 60 del artículo de Russell. No recomiendo leer en su totalidad el artículo de Strawson.

\paragraph{Para el artículo de Frege,} sugiero las siguientes palabras: relación, igualdad, conocimiento efectivo, sentido, referencia, nombre propio, sentido indirecto, representación, conidcional subjuntivo, condición \textit{salva veritate}, lenguaje lógicamente perfecto

\paragraph{Para el artículo de Russell,} sugiero las siguientes palabras: función descriptiva, descripción definida, descripción indefinida, funciones descriptivas, forma gramatical, Meinong, definición de la palabra `el (o la)', reconstruir el argumento que aparece al final de 51 (Es un párrafo), \textit{es} de predicación, \textit{es} de identidad, discutir en clase las pp. 57-60

\paragraph{Para el artículo de Strawson,} sugiero las siguientes palabras: uso referencial singularizador,

% \section{Notas del artículo de Strawson}
%Pienso que es correcto decir que la teoría de las descripciones de
%Rus- sell, que trata de la última de las cuatro clases de expresiones
%que he men- cionado más arriba (esto es, de expresiones de la forma
%«el tal-y-tal»), es aún ampliamente aceptada entre los lógicos por
%entender que proporcio- na una explicación correcta del uso de tales
%expresiones en el lenguaje ordi- nario. Deseo mostrar, en primer
%lugar, que esta teoría, entendida de tal mane- ra, contiene algunos
%errores fundamentales.
%\section{Notas del artículo de Russell}
%
% Hemos tenido ya ocasión de mencionar «funciones descriptivas», esto
% es: expresiones tales como «el padre de x» o «el seno de x». Estas
% expresiones han de definirse definiendo primero «descripciones». Una
% «descripción» puede ser de dos clases, definida o indefinida (o
% ambigua). Una descripción indefinida es una expresión de la forma
% «un tal-y-tal», y una descripción definida es una expresión de la
% forma «el tal-y-tal».
%
% \textbf{Glosario:} descripción definida, descripción indefinida,
% funciones descriptivas, forma gramatical, Meinong,
%
%\section{Notas del artículo de Frege}
%
%Esto surge de una imperfección del lenguaje, de la que tampoco está
%completamente libre por lo demás el lenguaje simbólico del análisi;
%también aquí pueden aparecer combinaciones de signos que producen la
%apariencia de que se refieren a algo pero que, por lo menos hasta
%ahora, carecen de referencia, como, por ejemplo, las series
%divergentes infinitas. Esto puede evitarse por medio de una
%estipulación especial al efecto de que, por ejemplo, las series
%divergentes infinitas deban referirse al número 0. De un lenguaje
%lógicamente perfecto (conceptografía) ha de reclamarse que cada
%expresión que se ha formado como nombre propio de modo
%gramaticalmente correcto, a partir de signos ya introducidos, designe
%también de hecho un objeto, y que no se introduzca ningún nuevo
%signo como nombre propio, sin que tenga asegurada una referencia. En
%los textos de lógica se advierte sobre la ambigiedad de las
%expresiones como fuente de errores lógicos.
%
%
%En el caso de estas oraciones, son posibles, por lo demás,
%concepciones ligeramente distintas. El sentido de la oración `
% `Después%de que Schleswig-Holstein se hubo separado de Dinamarca,
% Prusia y
%Austria se enemistaron'', puede reproducirse de la forma: ``Después
%de la separación de Schleswig-Holstein de Dinamarca, Prusia y Austria
%se enemistarom''. De acuerdo con esta concepción, resulta en verdad
%suficientemente claro que el pensamiento de que Schleswig-Holstein se
%separó alguna vez de Dinamarca no ha de tomarse como parte de este
%sentido, sino que ese pensamiento es la presuposición necesaria para
%que la expresión ``después de la separación de Schleswig-Holstein de
%Dinamarca'' tenga alguna referencia. Desde luego, nuestra oración
%puede concebirse de tal manera que se esté diciendo con ella que
%Schleswig-Holstein se separó alguna vez de Dinamarca. Tenemos
%entonces un caso que ha de considerarse más tarde. Para darse cuenta
%más claramente de la diferencia, pongámonos en el pellejo de un chino
%que, debido a su escaso conocimiento de la historia europea, tiene
%por falso el que Schleswig-Holstein se haya separado alguna vez de
%Dinamarca. Éste no tendría a nuestra oración, concebida del primer
%modo, ni por verdadera ni por falsa, sino que le negaría cualquier
%referencia, puesto que la oración subordinada carece de ella. Esta
%última sólo daría, aparentemente, una determinación temporal. Si, por
%el contrario, concibiese nuestra oración de acuerdo con el segundo
%modo, encontraría expresado en ella un pensamiento que tendría por
%falso, junto a una parte que, para él, carecería de referencia. %
%Cuando ha de indicarse un instante de manera indeterminada en el
%antecedente y el consecuente de una oración condicional, esto ocurre
%frecuentemente por medio del mero uso del tiempo presente del verbo
%que, en tal caso, no conlleva la designación del presente temporal.
%Esta forma gramatical es entonces el componente que indica de manera
%indeterminada en la oración principal y en la subordinada. ``Cuando
%el Sol está en el trópico de Cáncer, tenemos el día más largo en el
%hemisferio norte'' es un ejemplo de esto. Aquí es también imposible
%expresar el sentido de la subordinada en una principal, puesto que
%este sentido no es ningún pensamiento completo; pues si dijésemos:
%``El Sol está en el trópico de Cáncer'', estaríamos designando
% nuestro presente y, con ello, alteraríamos el sentido. Aún menos es
% el sentido de la principal un pensamiento; sólo el todo compuesto de
% la oración principal y de la subordinada contiene tal pensamiento.
% Por lo demás, también se pueden indicar de manera indeterminada
% diversos componentes comunes del antecedente y el consecuente de una
% oración condicional. Si encontramos que, por lo general, el valor
%cognoscitivo de ``a = a'' y ``a = b'' es diferente, esto se explica
%diciendo que, para el valor cognoscitivo, el sentido de la oración, a
%saber: el pensamiento %expresado por ella, no viene menos al caso que
%su referencia, esto es: su valor de verdad. Ahora bien, si a = b,
%entonces la referencia de `b' es ciertamente la misma que la de `a'
%y, por consiguiente, también el valor de verdad de `a = b' es el
%mismo que el de `a = a'. Con todo, el sentido de `b' puede ser
%distinto del sentido de `a', y con ello también el pensamiento
%expresado en `a = b' será diferente del expresado por `a = a'; en
%ese caso, ambas oraciones tampoco tienen el mismo valor
%cognoscitivo. Si entendemos por `juicio', como lo hemos hecho más
%arriba, el avance del pensamiento a su valor de verdad, diremos
%también entonces que los juicios son diferentes. Se comienza a
%partir en temas en 34 se vuelca al sentido y refProbablemente, uno se
%decidiría por esta opción. En el otro caso la situación estaría
%bastante enredada: tendríamos entonces más pensamientos simples que
%oraciones. Si sustituimos también ahora la oración «Napoleón cayó en
%la cuenta del peligro para su flanco derecho» por otra con el mismo
%valor de verdad, por ejemplo, por «Napoleón tenía ya más de cuarenta
%y cinco años», se habría alterado con ello no sólo nuestro primer
%pensamiento, sino tam- bién nuestro tercero, y por ello podría
%cambiar su valor de verdad, es decir: si su edad no era la razón de
%la decisión de dirigir su guardia contra el enemigo. Á partir de
%esto, puede percibirse por qué no siempre se pueden reemplazar
%entre sí las oraciones con el mismo valor de verdad. La oración
%expresa entonces, en virtud de su conexión con otra, más que por sí
%sola. Consideremos ahora casos donde tal cosa sucede regularmente. %En la oración «Bebel se imagina que la devolución de Alsacia-Lorena
%aplacará los deseos de venganza de Francia» erencia de oraciones, no
%de nombres. En la expresión del primer pensamiento, las palabras
%de la oración subordinada tienen su referencia indirecta, mientras
%que las mismas palabras tienen su referencia habitual en la expresión
%del segundo pensamiento. Vemos a partir de esto que la subordinada de
%nuestra oración compuesta originaria ha de tomarse ciertamente como
%doble con diferentes referencias, de las cuales una es un %pensamiento, la otra un valor de verdad. Ahora bien, puesto que el
%valor de verdad no es la referencia total de la oración subordinada,
%no podemos simplemente reemplazar ésta por otra con el mismo valor de
%verdad. Sucede algo similar con expresiones como «saber», «darse
%cuenta», «se sabe». Con una oración subordinada causal y su
%oración principal correspondiente, expresamos diversos pensamientos
%que, sin embargo, no corresponden a las oraciones tomadas de modo
%aislado. En la oración ``Puesto que el hielo es menos denso que el
%agua, flota en el agua'' tenemos: 1. El hielo es menos denso que el
%agua. 2. Si algo es menos denso que el agua, flota en el agua. 3. El
%hielo flota en el agua. El tercer pensamiento no necesitaba quizás
%ser introducido explícitamente, puesto que está contenido en los dos
%primeros. En cambio, ni combinando el primero con el tercero ni el
%segundo con el tercero, se lograría el sentido de nuestra oración.
%Puede verse ahora que en nuestra oración subordinada «puesto que el
%hielo es más denso que el agua» se expresa tanto nuestro primer
%pensamiento como una parte de nuestro segundo. De aquí viene el que
%no podamos reemplazar sin más nuestra subordinada por otra con el
%mismo valor de verdad; pues con ello se alteraría también nuestro
%segundo pensamiento y podría fácilmente verse afectado también su
%valor de verdad. El asunto es similar en la oración «Si el hierro
%fuera menos denso que el agua, flotaría en el agua». Tenemos aquí %los dos pensamientos siguientes: que el hierro no es menos denso que %el agua, y que algo flota en el %agua si es menos denso que el agua.


%En 36 es donde aparece el prototipo de la teoría deflacionista de la
%verdad y al final aparece el criterio \textit{salva veritate}

%37 oraciones subordinadas parte de actitudes proposicionales
%(contextos opacos)

%39 referencia indirecta, uso de órdenes (actos de habla)

%40 Descripciones definidas.

%21 Gottlob Frege, “On Sense and Nominatum”, (pp. 35-) Frege, “On Sense and Reference” (pp. 217-229)


%Russell, “On Denoting” (pp. 230-238)
%
% PReguntas para cuestionario de Sobre Sentido y REferencia
